% --- Archon physics formalization source begin ---
% archon:physics
% archon:covers PhyXMiniProblems/problem_phyx_mini_0665.lean
% archon:source-report reports/phyx_mini/problem_phyx_mini_0665.source.json

\chapter{Physics problem phyx\_mini\_0665}
\label{ch:PhyXMiniProblems_problem_phyx_mini_0665}

\paragraph{Problem source.}
\#\# Physical scenario

Liquid water with density \$\textbackslash{}rho\$ is filled on top of a thin piston in a cylinder with cross-sectional area \$A\$ and total height \$H\$, as shown in figure. Air is let in under the piston so that it pushes up, causing the water to spill over the edge.

\#\# Question

Derive the formula for the air pressure as a function of piston elevation from the bottom, \$h\$.

\#\# Answer choices

- A: A: P\_0 + (H + h)\textbackslash{}rho g

- B: B: P\_0 - (H - h)\textbackslash{}rho g

- C: C: P\_0 + \textbackslash{}frac\{1\}\{2\}(H - h)\textbackslash{}rho g

- D: D: P\_0 + (H - h)\textbackslash{}rho g

\#\# Figure caption (auxiliary; use the image as primary evidence)

The image depicts a diagram of a closed container filled partially with a fluid, illustrated in blue. The container is rectangular and has a thicker outer wall indicating insulation or structural support.

Key elements in the diagram:

1. **Fluid**: The fluid inside the container is shown as light blue with horizontal lines suggesting a liquid state.

2. **Height Indications**: 
   - There are two labeled heights on the left side: `H` and a smaller section labeled `h`, indicating different levels within the container.
   - `H` represents the full height of the liquid column in the container.
   - `h` might indicate a submerged or internal feature's height, but specific details aren't provided.

3. **Force and Direction**:
   - A downward arrow labeled with `g` represents gravitational force acting on the fluid.

4. **Air Connection**:
   - A line runs from the bottom section of the container to a symbol representing a valve, connected to "Air". This suggests an air pressure or release mechanism affecting the system.

5. **Structural Details**:
   - The container is shown with a thicker boundary, implying outer walls or insulation.

Overall, the image depicts a simplified cross-sectional view of a fluid-containing system, potentially for analyzing fluid dynamics, pressure, or buoyancy effects.

\#\# Dataset metadata

Dynamics; Physical Model Grounding Reasoning, Multi-Formula Reasoning

\paragraph{Recorded answer/context.}
D: D: P\_0 + (H - h)\textbackslash{}rho g

\paragraph{Figure/image path.}
/root/proposal\_for\_physic/science-mango/phyx\_mini\_run/phyx\_data/test\_image/665.png

\paragraph{Formalization target.}
create a compiling Lean file with sorry bodies at `PhyXMiniProblems/problem\_phyx\_mini\_0665.lean`. The Lean declarations must preserve the physical quantities, dimensions or dimensional roles, figure labels, governing-law hypotheses, and final relation expressed by this problem.
Use Mathlib/Physlib names found through LeanExplore where available. If a physics API is missing, introduce faithful local abstractions rather than scalar placeholder aliases.

\begin{theorem}[Physics formalization target]
\label{thm:physics:phyx_mini_0665:target}
\lean{PhyXMiniProblems.ProblemPhyXMini0665.problem_phyx_mini_0665}
\uses{def:physics:phyx-mini-0665:phyxminiproblems-problemphyxmini0665-lengthquantity, def:physics:phyx-mini-0665:phyxminiproblems-problemphyxmini0665-lengthreadout, def:physics:phyx-mini-0665:phyxminiproblems-problemphyxmini0665-densityreadout, def:physics:phyx-mini-0665:phyxminiproblems-problemphyxmini0665-accelerationreadout, def:physics:phyx-mini-0665:phyxminiproblems-problemphyxmini0665-pressurereadout, def:physics:phyx-mini-0665:phyxminiproblems-problemphyxmini0665-spillingwaterpistonsetup, def:physics:phyx-mini-0665:phyxminiproblems-problemphyxmini0665-matchesspillingwaterscenario, def:physics:phyx-mini-0665:phyxminiproblems-problemphyxmini0665-matchessuppliedpistoncylinderfigure, def:physics:phyx-mini-0665:phyxminiproblems-problemphyxmini0665-isadmissiblepistonelevation, def:physics:phyx-mini-0665:phyxminiproblems-problemphyxmini0665-hasphysicalspillingwaterparameters, def:physics:phyx-mini-0665:phyxminiproblems-problemphyxmini0665-satisfiesspillingcolumngeometry, def:physics:phyx-mini-0665:phyxminiproblems-problemphyxmini0665-satisfieshydrostaticpressurelaw, def:physics:phyx-mini-0665:phyxminiproblems-problemphyxmini0665-satisfiesthinpistonforcebalance}
The assigned autoformalize agent should translate this physics problem into Lean declarations in the covered file, with theorem and lemma proof bodies written as `by sorry`.
\end{theorem}
\begin{proof}
This is an autoformalization task, not a proof task. Produce statements that can later be proved by the physics prover without weakening the physical model.
\end{proof}

% --- Archon declaration topology begin ---
\section{Declaration topology}
\begin{definition}[mass Density Dimension]
  \label{def:physics:phyx-mini-0665:phyxminiproblems-problemphyxmini0665-massdensitydimension}
  \lean{PhyXMiniProblems.ProblemPhyXMini0665.massDensityDimension}
  The physical dimension of mass density, M L\textasciicircum{}-3
\end{definition}
\begin{proof}
  This is the declared model definition of mass Density Dimension.
\end{proof}
\begin{definition}[acceleration Dimension]
  \label{def:physics:phyx-mini-0665:phyxminiproblems-problemphyxmini0665-accelerationdimension}
  \lean{PhyXMiniProblems.ProblemPhyXMini0665.accelerationDimension}
  The physical dimension of acceleration, L T\textasciicircum{}-2
\end{definition}
\begin{proof}
  This is the declared model definition of acceleration Dimension.
\end{proof}
\begin{definition}[Length Quantity]
  \label{def:physics:phyx-mini-0665:phyxminiproblems-problemphyxmini0665-lengthquantity}
  \lean{PhyXMiniProblems.ProblemPhyXMini0665.LengthQuantity}
  A nonnegative physical length, independent of the readout unit
\end{definition}
\begin{proof}
  This is the declared model definition of Length Quantity.
\end{proof}
\begin{definition}[Mass Density Quantity]
  \label{def:physics:phyx-mini-0665:phyxminiproblems-problemphyxmini0665-massdensityquantity}
  \lean{PhyXMiniProblems.ProblemPhyXMini0665.MassDensityQuantity}
  \uses{def:physics:phyx-mini-0665:phyxminiproblems-problemphyxmini0665-massdensitydimension}
  A nonnegative, spatially uniform mass density
\end{definition}
\begin{proof}
  This is the declared model definition of Mass Density Quantity.
\end{proof}
\begin{definition}[Acceleration Quantity]
  \label{def:physics:phyx-mini-0665:phyxminiproblems-problemphyxmini0665-accelerationquantity}
  \lean{PhyXMiniProblems.ProblemPhyXMini0665.AccelerationQuantity}
  \uses{def:physics:phyx-mini-0665:phyxminiproblems-problemphyxmini0665-accelerationdimension}
  A nonnegative acceleration magnitude
\end{definition}
\begin{proof}
  This is the declared model definition of Acceleration Quantity.
\end{proof}
\begin{definition}[Pressure Quantity]
  \label{def:physics:phyx-mini-0665:phyxminiproblems-problemphyxmini0665-pressurequantity}
  \lean{PhyXMiniProblems.ProblemPhyXMini0665.PressureQuantity}
  A physical pressure, using Physlib's dimension M L\textasciicircum{}-1 T\textasciicircum{}-2
\end{definition}
\begin{proof}
  This is the declared model definition of Pressure Quantity.
\end{proof}
\begin{definition}[length Readout]
  \label{def:physics:phyx-mini-0665:phyxminiproblems-problemphyxmini0665-lengthreadout}
  \lean{PhyXMiniProblems.ProblemPhyXMini0665.lengthReadout}
  \uses{def:physics:phyx-mini-0665:phyxminiproblems-problemphyxmini0665-lengthquantity}
  Read a physical length in the length unit selected by units
\end{definition}
\begin{proof}
  This is the declared model definition of length Readout.
\end{proof}
\begin{definition}[area Readout]
  \label{def:physics:phyx-mini-0665:phyxminiproblems-problemphyxmini0665-areareadout}
  \lean{PhyXMiniProblems.ProblemPhyXMini0665.areaReadout}
  Read an area in the squared length unit selected by units
\end{definition}
\begin{proof}
  This is the declared model definition of area Readout.
\end{proof}
\begin{definition}[density Readout]
  \label{def:physics:phyx-mini-0665:phyxminiproblems-problemphyxmini0665-densityreadout}
  \lean{PhyXMiniProblems.ProblemPhyXMini0665.densityReadout}
  \uses{def:physics:phyx-mini-0665:phyxminiproblems-problemphyxmini0665-massdensityquantity}
  Read a mass density in the coherent units selected by units
\end{definition}
\begin{proof}
  This is the declared model definition of density Readout.
\end{proof}
\begin{definition}[acceleration Readout]
  \label{def:physics:phyx-mini-0665:phyxminiproblems-problemphyxmini0665-accelerationreadout}
  \lean{PhyXMiniProblems.ProblemPhyXMini0665.accelerationReadout}
  \uses{def:physics:phyx-mini-0665:phyxminiproblems-problemphyxmini0665-accelerationquantity}
  Read an acceleration magnitude in the coherent units selected by units
\end{definition}
\begin{proof}
  This is the declared model definition of acceleration Readout.
\end{proof}
\begin{definition}[pressure Readout]
  \label{def:physics:phyx-mini-0665:phyxminiproblems-problemphyxmini0665-pressurereadout}
  \lean{PhyXMiniProblems.ProblemPhyXMini0665.pressureReadout}
  \uses{def:physics:phyx-mini-0665:phyxminiproblems-problemphyxmini0665-pressurequantity}
  Read a pressure in the coherent pressure unit induced by units
\end{definition}
\begin{proof}
  This is the declared model definition of pressure Readout.
\end{proof}
\begin{definition}[length In Meters]
  \label{def:physics:phyx-mini-0665:phyxminiproblems-problemphyxmini0665-lengthinmeters}
  \lean{PhyXMiniProblems.ProblemPhyXMini0665.lengthInMeters}
  \uses{def:physics:phyx-mini-0665:phyxminiproblems-problemphyxmini0665-lengthquantity, def:physics:phyx-mini-0665:phyxminiproblems-problemphyxmini0665-lengthreadout}
  Metre readout of a physical length
\end{definition}
\begin{proof}
  This is the declared model definition of length In Meters.
\end{proof}
\begin{definition}[area In Square Meters]
  \label{def:physics:phyx-mini-0665:phyxminiproblems-problemphyxmini0665-areainsquaremeters}
  \lean{PhyXMiniProblems.ProblemPhyXMini0665.areaInSquareMeters}
  \uses{def:physics:phyx-mini-0665:phyxminiproblems-problemphyxmini0665-areareadout}
  Square-metre readout of a physical area
\end{definition}
\begin{proof}
  This is the declared model definition of area In Square Meters.
\end{proof}
\begin{definition}[density In Kilograms Per Cubic Meter]
  \label{def:physics:phyx-mini-0665:phyxminiproblems-problemphyxmini0665-densityinkilogramspercubicmeter}
  \lean{PhyXMiniProblems.ProblemPhyXMini0665.densityInKilogramsPerCubicMeter}
  \uses{def:physics:phyx-mini-0665:phyxminiproblems-problemphyxmini0665-massdensityquantity, def:physics:phyx-mini-0665:phyxminiproblems-problemphyxmini0665-densityreadout}
  Kilogram-per-cubic-metre readout of a physical mass density
\end{definition}
\begin{proof}
  This is the declared model definition of density In Kilograms Per Cubic Meter.
\end{proof}
\begin{definition}[acceleration In Meters Per Second Squared]
  \label{def:physics:phyx-mini-0665:phyxminiproblems-problemphyxmini0665-accelerationinmeterspersecondsquared}
  \lean{PhyXMiniProblems.ProblemPhyXMini0665.accelerationInMetersPerSecondSquared}
  \uses{def:physics:phyx-mini-0665:phyxminiproblems-problemphyxmini0665-accelerationquantity, def:physics:phyx-mini-0665:phyxminiproblems-problemphyxmini0665-accelerationreadout}
  Metre-per-second-squared readout of an acceleration magnitude
\end{definition}
\begin{proof}
  This is the declared model definition of acceleration In Meters Per Second Squared.
\end{proof}
\begin{definition}[pressure In Pascals]
  \label{def:physics:phyx-mini-0665:phyxminiproblems-problemphyxmini0665-pressureinpascals}
  \lean{PhyXMiniProblems.ProblemPhyXMini0665.pressureInPascals}
  \uses{def:physics:phyx-mini-0665:phyxminiproblems-problemphyxmini0665-pressurequantity, def:physics:phyx-mini-0665:phyxminiproblems-problemphyxmini0665-pressurereadout}
  Pascal readout of a physical pressure
\end{definition}
\begin{proof}
  This is the declared model definition of pressure In Pascals.
\end{proof}
\begin{definition}[Figure Label]
  \label{def:physics:phyx-mini-0665:phyxminiproblems-problemphyxmini0665-figurelabel}
  \lean{PhyXMiniProblems.ProblemPhyXMini0665.FigureLabel}
  Literal labels visible in the primary raster
\end{definition}
\begin{proof}
  This is the declared model definition of Figure Label.
\end{proof}
\begin{definition}[Height Mark]
  \label{def:physics:phyx-mini-0665:phyxminiproblems-problemphyxmini0665-heightmark}
  \lean{PhyXMiniProblems.ProblemPhyXMini0665.HeightMark}
  The two vertical dimension marks drawn to the left of the cylinder
\end{definition}
\begin{proof}
  This is the declared model definition of Height Mark.
\end{proof}
\begin{definition}[Figure Location]
  \label{def:physics:phyx-mini-0665:phyxminiproblems-problemphyxmini0665-figurelocation}
  \lean{PhyXMiniProblems.ProblemPhyXMini0665.FigureLocation}
  Distinguished vertical locations in the cross-sectional diagram
\end{definition}
\begin{proof}
  This is the declared model definition of Figure Location.
\end{proof}
\begin{definition}[Vertical Direction]
  \label{def:physics:phyx-mini-0665:phyxminiproblems-problemphyxmini0665-verticaldirection}
  \lean{PhyXMiniProblems.ProblemPhyXMini0665.VerticalDirection}
  Vertical directions used to interpret the arrow labelled g
\end{definition}
\begin{proof}
  This is the declared model definition of Vertical Direction.
\end{proof}
\begin{definition}[Cylinder Liquid]
  \label{def:physics:phyx-mini-0665:phyxminiproblems-problemphyxmini0665-cylinderliquid}
  \lean{PhyXMiniProblems.ProblemPhyXMini0665.CylinderLiquid}
  Working liquid named by the problem
\end{definition}
\begin{proof}
  This is the declared model definition of Cylinder Liquid.
\end{proof}
\begin{definition}[Supplied Gas]
  \label{def:physics:phyx-mini-0665:phyxminiproblems-problemphyxmini0665-suppliedgas}
  \lean{PhyXMiniProblems.ProblemPhyXMini0665.SuppliedGas}
  Gas supplied through the valve below the piston
\end{definition}
\begin{proof}
  This is the declared model definition of Supplied Gas.
\end{proof}
\begin{definition}[Piston Cylinder Figure]
  \label{def:physics:phyx-mini-0665:phyxminiproblems-problemphyxmini0665-pistoncylinderfigure}
  \lean{PhyXMiniProblems.ProblemPhyXMini0665.PistonCylinderFigure}
  \uses{def:physics:phyx-mini-0665:phyxminiproblems-problemphyxmini0665-lengthquantity, def:physics:phyx-mini-0665:phyxminiproblems-problemphyxmini0665-figurelabel, def:physics:phyx-mini-0665:phyxminiproblems-problemphyxmini0665-heightmark, def:physics:phyx-mini-0665:phyxminiproblems-problemphyxmini0665-figurelocation, def:physics:phyx-mini-0665:phyxminiproblems-problemphyxmini0665-verticaldirection}
  Structured transcription of the physical information visible in image 665.png. The label h denotes a variable piston elevation rather than a fixed numerical readout
\end{definition}
\begin{proof}
  This is the declared model definition of Piston Cylinder Figure.
\end{proof}
\begin{definition}[Spilling Water Piston Setup]
  \label{def:physics:phyx-mini-0665:phyxminiproblems-problemphyxmini0665-spillingwaterpistonsetup}
  \lean{PhyXMiniProblems.ProblemPhyXMini0665.SpillingWaterPistonSetup}
  \uses{def:physics:phyx-mini-0665:phyxminiproblems-problemphyxmini0665-lengthquantity, def:physics:phyx-mini-0665:phyxminiproblems-problemphyxmini0665-massdensityquantity, def:physics:phyx-mini-0665:phyxminiproblems-problemphyxmini0665-accelerationquantity, def:physics:phyx-mini-0665:phyxminiproblems-problemphyxmini0665-pressurequantity, def:physics:phyx-mini-0665:phyxminiproblems-problemphyxmini0665-cylinderliquid, def:physics:phyx-mini-0665:phyxminiproblems-problemphyxmini0665-suppliedgas, def:physics:phyx-mini-0665:phyxminiproblems-problemphyxmini0665-pistoncylinderfigure}
  The cylinder and its independent physical observables. In particular, airPressureUnderPiston, waterPressureOnPiston, and waterColumnDepthAbovePiston are fields rather than definitions of the requested formula
\end{definition}
\begin{proof}
  This is the declared model definition of Spilling Water Piston Setup.
\end{proof}
\begin{definition}[Matches Spilling Water Scenario]
  \label{def:physics:phyx-mini-0665:phyxminiproblems-problemphyxmini0665-matchesspillingwaterscenario}
  \lean{PhyXMiniProblems.ProblemPhyXMini0665.MatchesSpillingWaterScenario}
  \uses{def:physics:phyx-mini-0665:phyxminiproblems-problemphyxmini0665-spillingwaterpistonsetup}
  Qualitative information supplied by the problem prose
\end{definition}
\begin{proof}
  This is the declared model definition of Matches Spilling Water Scenario.
\end{proof}
\begin{definition}[Matches Supplied Piston Cylinder Figure]
  \label{def:physics:phyx-mini-0665:phyxminiproblems-problemphyxmini0665-matchessuppliedpistoncylinderfigure}
  \lean{PhyXMiniProblems.ProblemPhyXMini0665.MatchesSuppliedPistonCylinderFigure}
  \uses{def:physics:phyx-mini-0665:phyxminiproblems-problemphyxmini0665-spillingwaterpistonsetup}
  Literal evidence from the supplied raster. It contains no pressure formula and no answer-choice selection
\end{definition}
\begin{proof}
  This is the declared model definition of Matches Supplied Piston Cylinder Figure.
\end{proof}
\begin{definition}[Is Admissible Piston Elevation]
  \label{def:physics:phyx-mini-0665:phyxminiproblems-problemphyxmini0665-isadmissiblepistonelevation}
  \lean{PhyXMiniProblems.ProblemPhyXMini0665.IsAdmissiblePistonElevation}
  \uses{def:physics:phyx-mini-0665:phyxminiproblems-problemphyxmini0665-lengthquantity, def:physics:phyx-mini-0665:phyxminiproblems-problemphyxmini0665-lengthinmeters, def:physics:phyx-mini-0665:phyxminiproblems-problemphyxmini0665-spillingwaterpistonsetup}
  A piston elevation lies between the bottom datum and the top rim
\end{definition}
\begin{proof}
  This is the declared model definition of Is Admissible Piston Elevation.
\end{proof}
\begin{definition}[Has Physical Spilling Water Parameters]
  \label{def:physics:phyx-mini-0665:phyxminiproblems-problemphyxmini0665-hasphysicalspillingwaterparameters}
  \lean{PhyXMiniProblems.ProblemPhyXMini0665.HasPhysicalSpillingWaterParameters}
  \uses{def:physics:phyx-mini-0665:phyxminiproblems-problemphyxmini0665-lengthinmeters, def:physics:phyx-mini-0665:phyxminiproblems-problemphyxmini0665-areainsquaremeters, def:physics:phyx-mini-0665:phyxminiproblems-problemphyxmini0665-densityinkilogramspercubicmeter, def:physics:phyx-mini-0665:phyxminiproblems-problemphyxmini0665-accelerationinmeterspersecondsquared, def:physics:phyx-mini-0665:phyxminiproblems-problemphyxmini0665-pressureinpascals, def:physics:phyx-mini-0665:phyxminiproblems-problemphyxmini0665-spillingwaterpistonsetup}
  Positivity conditions needed for a physically meaningful force balance
\end{definition}
\begin{proof}
  This is the declared model definition of Has Physical Spilling Water Parameters.
\end{proof}
\begin{definition}[Satisfies Spilling Column Geometry]
  \label{def:physics:phyx-mini-0665:phyxminiproblems-problemphyxmini0665-satisfiesspillingcolumngeometry}
  \lean{PhyXMiniProblems.ProblemPhyXMini0665.SatisfiesSpillingColumnGeometry}
  \uses{def:physics:phyx-mini-0665:phyxminiproblems-problemphyxmini0665-lengthquantity, def:physics:phyx-mini-0665:phyxminiproblems-problemphyxmini0665-lengthreadout, def:physics:phyx-mini-0665:phyxminiproblems-problemphyxmini0665-spillingwaterpistonsetup, def:physics:phyx-mini-0665:phyxminiproblems-problemphyxmini0665-isadmissiblepistonelevation}
  Because water spills over the open rim, the free-surface elevation remains H; hence the water depth above a piston at elevation h is H - h
\end{definition}
\begin{proof}
  This is the declared model definition of Satisfies Spilling Column Geometry.
\end{proof}
\begin{definition}[Satisfies Hydrostatic Pressure Law]
  \label{def:physics:phyx-mini-0665:phyxminiproblems-problemphyxmini0665-satisfieshydrostaticpressurelaw}
  \lean{PhyXMiniProblems.ProblemPhyXMini0665.SatisfiesHydrostaticPressureLaw}
  \uses{def:physics:phyx-mini-0665:phyxminiproblems-problemphyxmini0665-lengthquantity, def:physics:phyx-mini-0665:phyxminiproblems-problemphyxmini0665-lengthreadout, def:physics:phyx-mini-0665:phyxminiproblems-problemphyxmini0665-densityreadout, def:physics:phyx-mini-0665:phyxminiproblems-problemphyxmini0665-accelerationreadout, def:physics:phyx-mini-0665:phyxminiproblems-problemphyxmini0665-pressurereadout, def:physics:phyx-mini-0665:phyxminiproblems-problemphyxmini0665-spillingwaterpistonsetup, def:physics:phyx-mini-0665:phyxminiproblems-problemphyxmini0665-isadmissiblepistonelevation}
  Hydrostatic equilibrium between the ambient open free surface and the water face of the piston. This law uses the independent water-column depth; it does not state the requested air-pressure formula
\end{definition}
\begin{proof}
  This is the declared model definition of Satisfies Hydrostatic Pressure Law.
\end{proof}
\begin{definition}[Satisfies Thin Piston Force Balance]
  \label{def:physics:phyx-mini-0665:phyxminiproblems-problemphyxmini0665-satisfiesthinpistonforcebalance}
  \lean{PhyXMiniProblems.ProblemPhyXMini0665.SatisfiesThinPistonForceBalance}
  \uses{def:physics:phyx-mini-0665:phyxminiproblems-problemphyxmini0665-lengthquantity, def:physics:phyx-mini-0665:phyxminiproblems-problemphyxmini0665-areareadout, def:physics:phyx-mini-0665:phyxminiproblems-problemphyxmini0665-pressurereadout, def:physics:phyx-mini-0665:phyxminiproblems-problemphyxmini0665-spillingwaterpistonsetup, def:physics:phyx-mini-0665:phyxminiproblems-problemphyxmini0665-isadmissiblepistonelevation}
  For the thin, weightless, frictionless piston in quasistatic balance, the air force upward equals the water-pressure force downward. The common area A is retained explicitly rather than cancelled inside the assumption
\end{definition}
\begin{proof}
  This is the declared model definition of Satisfies Thin Piston Force Balance.
\end{proof}
\begin{definition}[Answer Choice]
  \label{def:physics:phyx-mini-0665:phyxminiproblems-problemphyxmini0665-answerchoice}
  \lean{PhyXMiniProblems.ProblemPhyXMini0665.AnswerChoice}
  Labels of the four formulas printed in the problem
\end{definition}
\begin{proof}
  This is the declared model definition of Answer Choice.
\end{proof}
\begin{definition}[displayed Pressure Formula]
  \label{def:physics:phyx-mini-0665:phyxminiproblems-problemphyxmini0665-displayedpressureformula}
  \lean{PhyXMiniProblems.ProblemPhyXMini0665.displayedPressureFormula}
  \uses{def:physics:phyx-mini-0665:phyxminiproblems-problemphyxmini0665-lengthquantity, def:physics:phyx-mini-0665:phyxminiproblems-problemphyxmini0665-lengthreadout, def:physics:phyx-mini-0665:phyxminiproblems-problemphyxmini0665-densityreadout, def:physics:phyx-mini-0665:phyxminiproblems-problemphyxmini0665-accelerationreadout, def:physics:phyx-mini-0665:phyxminiproblems-problemphyxmini0665-pressurereadout, def:physics:phyx-mini-0665:phyxminiproblems-problemphyxmini0665-spillingwaterpistonsetup, def:physics:phyx-mini-0665:phyxminiproblems-problemphyxmini0665-answerchoice}
  The pressure formula printed beside each answer label, evaluated in one coherent unit system. This is presentation data only; the main theorem below states the physical target explicitly and does not unfold this table
\end{definition}
\begin{proof}
  This is the declared model definition of displayed Pressure Formula.
\end{proof}
\begin{definition}[Is Correct Answer]
  \label{def:physics:phyx-mini-0665:phyxminiproblems-problemphyxmini0665-iscorrectanswer}
  \lean{PhyXMiniProblems.ProblemPhyXMini0665.IsCorrectAnswer}
  \uses{def:physics:phyx-mini-0665:phyxminiproblems-problemphyxmini0665-lengthquantity, def:physics:phyx-mini-0665:phyxminiproblems-problemphyxmini0665-pressurereadout, def:physics:phyx-mini-0665:phyxminiproblems-problemphyxmini0665-spillingwaterpistonsetup, def:physics:phyx-mini-0665:phyxminiproblems-problemphyxmini0665-answerchoice, def:physics:phyx-mini-0665:phyxminiproblems-problemphyxmini0665-displayedpressureformula}
  A printed choice agrees with the physical air pressure in every coherent unit system
\end{definition}
\begin{proof}
  This is the declared model definition of Is Correct Answer.
\end{proof}
\begin{lemma}[answer Choice D is Correct]
  \label{lem:physics:phyx-mini-0665:phyxminiproblems-problemphyxmini0665-answerchoiced-iscorrect}
  \lean{PhyXMiniProblems.ProblemPhyXMini0665.answerChoiceD_isCorrect}
  \uses{def:physics:phyx-mini-0665:phyxminiproblems-problemphyxmini0665-lengthquantity, def:physics:phyx-mini-0665:phyxminiproblems-problemphyxmini0665-spillingwaterpistonsetup, def:physics:phyx-mini-0665:phyxminiproblems-problemphyxmini0665-matchesspillingwaterscenario, def:physics:phyx-mini-0665:phyxminiproblems-problemphyxmini0665-matchessuppliedpistoncylinderfigure, def:physics:phyx-mini-0665:phyxminiproblems-problemphyxmini0665-isadmissiblepistonelevation, def:physics:phyx-mini-0665:phyxminiproblems-problemphyxmini0665-hasphysicalspillingwaterparameters, def:physics:phyx-mini-0665:phyxminiproblems-problemphyxmini0665-satisfiesspillingcolumngeometry, def:physics:phyx-mini-0665:phyxminiproblems-problemphyxmini0665-satisfieshydrostaticpressurelaw, def:physics:phyx-mini-0665:phyxminiproblems-problemphyxmini0665-satisfiesthinpistonforcebalance, def:physics:phyx-mini-0665:phyxminiproblems-problemphyxmini0665-iscorrectanswer}
  The explicit pressure formula selects answer choice D
\end{lemma}
\begin{proof}
  The conclusion follows from the governing relations and figure readouts named in the statement of answer Choice D is Correct.
\end{proof}
% --- Archon declaration topology end ---
% --- Archon physics formalization source end ---
